%==========================================================
%  RAPPORT DE STAGE PFA — main.tex
%  Document principal — inclut tous les chapitres
%
%  COMPILATION :
%    pdflatex main.tex
%    bibtex   main
%    pdflatex main.tex
%    pdflatex main.tex
%==========================================================
\documentclass[a4paper,12pt,oneside]{report}

%──────────────────────────────────────────────────────────
%  ENCODAGE & LANGUE
%──────────────────────────────────────────────────────────
\usepackage[utf8]{inputenc}
\usepackage[T1]{fontenc}
\usepackage[french,english]{babel}   % français par défaut, anglais pour abstract
\usepackage{lmodern}
\usepackage{microtype}

%──────────────────────────────────────────────────────────
%  MISE EN PAGE
%──────────────────────────────────────────────────────────
\usepackage[
  top=2.5cm, bottom=2.8cm,
  left=3cm, right=2.5cm,
  headheight=14pt
]{geometry}
\usepackage{fancyhdr}
\usepackage{emptypage}
\usepackage{setspace}
\onehalfspacing

%──────────────────────────────────────────────────────────
%  COULEURS
%──────────────────────────────────────────────────────────
\usepackage[dvipsnames,table]{xcolor}
\definecolor{ChapBlue}{RGB}{0,51,102}
\definecolor{SecGrey}{RGB}{80,80,95}
\definecolor{AccentBlue}{RGB}{30,120,190}
\definecolor{LightBg}{RGB}{240,244,250}
\definecolor{CodeBg}{RGB}{248,248,250}
\definecolor{WarningOrange}{RGB}{220,120,0}

%──────────────────────────────────────────────────────────
%  GRAPHIQUES & FIGURES
%──────────────────────────────────────────────────────────
\usepackage{graphicx}
\usepackage{float}
\usepackage{subfigure}
\usepackage{tikz}
\usetikzlibrary{shapes.geometric, arrows.meta, positioning, calc, shadows}
\graphicspath{{./figures/}{./chapters/}}

%──────────────────────────────────────────────────────────
%  TABLEAUX
%──────────────────────────────────────────────────────────
\usepackage{booktabs}
\usepackage{array}
\usepackage{tabularx}
\usepackage{multirow}
\usepackage{longtable}
\renewcommand{\arraystretch}{1.3}

%──────────────────────────────────────────────────────────
%  LISTES
%──────────────────────────────────────────────────────────
\usepackage{enumitem}
\setlist[itemize,1]{label=\textcolor{AccentBlue}{\textbullet}, leftmargin=1.5em}
\setlist[itemize,2]{label=\textcolor{SecGrey}{\guillemotright}, leftmargin=2em}
\setlist[enumerate,1]{leftmargin=1.8em, label=\textcolor{AccentBlue}{\arabic*.}}

%──────────────────────────────────────────────────────────
%  CODE SOURCE
%──────────────────────────────────────────────────────────
\usepackage{listings}
\lstdefinestyle{PFAStyle}{
  backgroundcolor=\color{CodeBg},
  basicstyle=\ttfamily\small,
  keywordstyle=\color{ChapBlue}\bfseries,
  commentstyle=\color{SecGrey}\itshape,
  stringstyle=\color{AccentBlue},
  numberstyle=\tiny\color{SecGrey},
  numbers=left, numbersep=8pt,
  stepnumber=1,
  breaklines=true,
  breakatwhitespace=false,
  tabsize=2,
  frame=single,
  framesep=4pt,
  framerule=0.4pt,
  rulecolor=\color{AccentBlue!40},
  showspaces=false,
  showstringspaces=false,
  captionpos=b,
  xleftmargin=15pt
}
\lstset{style=PFAStyle}

% Langages supplémentaires
\lstdefinelanguage{JSON}{
  morestring=[b]",
  morecomment=[l]{//},
  morekeywords={true,false,null}
}

%──────────────────────────────────────────────────────────
%  STYLE DES CHAPITRES & SECTIONS
%──────────────────────────────────────────────────────────
\usepackage{titlesec}
\usepackage{titletoc}

% Format du titre de chapitre
\titleformat{\chapter}[display]
  {\normalfont\sffamily\huge\bfseries\color{ChapBlue}}
  {\tikz[baseline=(n.base)] \node[
      fill=ChapBlue, text=white, rounded corners=4pt,
      inner sep=8pt, font=\sffamily\Large\bfseries
    ] (n) {\chaptertitlename\ \thechapter};}
  {10pt}
  {\Huge}
  [\vspace{4pt}\color{AccentBlue}\titlerule[1.5pt]\vspace{2pt}
   \color{AccentBlue!40}\titlerule[0.4pt]]

\titlespacing*{\chapter}{0pt}{0pt}{30pt}

\titleformat{\section}
  {\color{ChapBlue}\large\bfseries\sffamily}
  {\thesection.}{0.6em}{}
  [\color{AccentBlue}\titlerule[0.6pt]]

\titleformat{\subsection}
  {\color{SecGrey}\normalsize\bfseries\sffamily}
  {\thesubsection.}{0.5em}{}

\titleformat{\subsubsection}
  {\color{SecGrey}\small\bfseries\itshape}
  {\thesubsubsection.}{0.4em}{}

%──────────────────────────────────────────────────────────
%  EN-TÊTES / PIEDS DE PAGE
%──────────────────────────────────────────────────────────
\pagestyle{fancy}
\fancyhf{}
\renewcommand{\headrulewidth}{0.4pt}
\renewcommand{\footrulewidth}{0.4pt}
\fancyhead[L]{\small\color{SecGrey}\sffamily\leftmark}
\fancyhead[R]{\small\color{SecGrey}\sffamily Rapport de Stage PFA}
\fancyfoot[C]{\small\color{SecGrey}\thepage}
\fancyfoot[L]{\small\color{SecGrey}\sffamily [À COMPLÉTER — Nom Étudiant]}
\fancyfoot[R]{\small\color{SecGrey}\sffamily \AcademicYear}

% Style pour les premières pages de chapitre
\fancypagestyle{plain}{
  \fancyhf{}
  \renewcommand{\headrulewidth}{0pt}
  \renewcommand{\footrulewidth}{0.4pt}
  \fancyfoot[C]{\small\color{SecGrey}\thepage}
}

%──────────────────────────────────────────────────────────
%  BOÎTES (tcolorbox)
%──────────────────────────────────────────────────────────
\usepackage[most]{tcolorbox}
\tcbuselibrary{skins,breakable,listings}

\newtcolorbox{remarque}[1][]{
  enhanced, breakable,
  title=Remarque,
  colback=LightBg, colframe=AccentBlue,
  fonttitle=\sffamily\bfseries\small,
  attach boxed title to top left={xshift=6pt,yshift=-3pt},
  boxed title style={colback=AccentBlue},
  #1
}

\newtcolorbox{definition}[2][]{
  enhanced, breakable,
  title=Définition : #2,
  colback=white, colframe=ChapBlue,
  fonttitle=\sffamily\bfseries\small,
  attach boxed title to top left={xshift=6pt,yshift=-3pt},
  boxed title style={colback=ChapBlue},
  #1
}

\newtcolorbox{important}[1][]{
  enhanced, breakable,
  colback=WarningOrange!8, colframe=WarningOrange,
  leftrule=4pt, toprule=0.5pt, bottomrule=0.5pt, rightrule=0.5pt,
  fontupper=\small,
  #1
}

%──────────────────────────────────────────────────────────
%  MATHÉMATIQUES
%──────────────────────────────────────────────────────────
\usepackage{amsmath,amssymb,amsfonts}

%──────────────────────────────────────────────────────────
%  HYPERLIENS (toujours en dernier ou avant geometry)
%──────────────────────────────────────────────────────────
\usepackage[
  colorlinks=true,
  linkcolor=ChapBlue,
  urlcolor=AccentBlue,
  citecolor=SecGrey,
  filecolor=ChapBlue,
  pdfauthor={[À COMPLÉTER — Nom Étudiant]},
  pdftitle={Rapport de Stage PFA — [À COMPLÉTER]},
  pdfsubject={Rapport de Stage de Fin d'Année},
  pdfkeywords={PFA, stage, rapport, [À COMPLÉTER]},
  pdfpagemode=UseOutlines,
  bookmarks=true,
  bookmarksnumbered=true,
  bookmarksopen=true
]{hyperref}
\usepackage{bookmark}

%──────────────────────────────────────────────────────────
%  BIBLIOGRAPHIE
%──────────────────────────────────────────────────────────
\usepackage{cite}   % style IEEE (numérotation [1])
% Pour style APA, remplacer par : \usepackage[authoryear]{natbib}

%──────────────────────────────────────────────────────────
%  COMMANDES UTILITAIRES
%──────────────────────────────────────────────────────────
\newcommand{\todo}[1]{\textcolor{red}{\textbf{[À COMPLÉTER : #1]}}}
\newcommand{\placeholder}[1]{%
  \begin{center}
    \fcolorbox{AccentBlue!50}{LightBg}{%
      \parbox{0.85\textwidth}{\centering\color{SecGrey}\itshape
        \textbf{[Placeholder]} #1}}
  \end{center}}

% Tag coloré inline
\newcommand{\tech}[1]{%
  \tikz[baseline=(t.base)]
    \node[rounded corners=3pt,fill=AccentBlue!12,draw=AccentBlue!40,
          inner sep=2pt 3pt,font=\scriptsize\sffamily\bfseries,
          text=ChapBlue](t){#1};}

% Raccourci sigle
\newcommand{\sigla}[2]{\textbf{#1} (#2)}

%──────────────────────────────────────────────────────────
%  MÉTADONNÉES STAGE (à personnaliser)
%──────────────────────────────────────────────────────────
\newcommand{\StudentName}{[À COMPLÉTER — Prénom NOM]}
\newcommand{\StudentField}{[À COMPLÉTER — Filière]}
\newcommand{\StudentYear}{[À COMPLÉTER — ex: 3\textsuperscript{e} année]}
\newcommand{\SchoolName}{[À COMPLÉTER — École / Université]}
\newcommand{\CompanyName}{[À COMPLÉTER — Entreprise]}
\newcommand{\CompanyCity}{[À COMPLÉTER — Ville]}
\newcommand{\ProjectTitle}{[À COMPLÉTER — Titre du Projet]}
\newcommand{\SupervisorPro}{[À COMPLÉTER — Encadrant Professionnel]}
\newcommand{\SupervisorAca}{[À COMPLÉTER — Encadrant Académique]}
\newcommand{\InternshipStart}{[À COMPLÉTER — Date début]}
\newcommand{\InternshipEnd}{[À COMPLÉTER — Date fin]}
\newcommand{\AcademicYear}{2024--2025}
\newcommand{\SubmissionDate}{[À COMPLÉTER — Mois Année]}

%==========================================================
%  DOCUMENT
%==========================================================
\begin{document}
\pagenumbering{gobble}   % pas de numéros sur la page de garde

%──────────────────────────────────────────────────────────
%  PAGE DE GARDE
%──────────────────────────────────────────────────────────
%==========================================================
%  CHAPITRE 00 — PAGE DE GARDE OFFICIELLE
%  Inclus depuis main.tex via %==========================================================
%  CHAPITRE 00 — PAGE DE GARDE OFFICIELLE
%  Inclus depuis main.tex via %==========================================================
%  CHAPITRE 00 — PAGE DE GARDE OFFICIELLE
%  Inclus depuis main.tex via \include{chapters/00_titlepage}
%==========================================================
\begin{titlepage}
\thispagestyle{empty}

% Couche TikZ pour les éléments graphiques de fond
\begin{tikzpicture}[remember picture, overlay]
  % Bandeau supérieur bleu foncé
  \fill[ChapBlue]
    (current page.north west) rectangle
    ([yshift=-4.5cm] current page.north east);
  % Liseré accent
  \fill[AccentBlue]
    ([yshift=-4.5cm] current page.north west) rectangle
    ([yshift=-4.8cm] current page.north east);
  % Barre latérale gauche
  \fill[ChapBlue!15]
    (current page.south west) rectangle
    ([xshift=0.7cm] current page.north west);
  % Bandeau bas
  \fill[ChapBlue]
    (current page.south west) rectangle
    ([yshift=2cm] current page.south east);
  \fill[AccentBlue]
    ([yshift=2cm] current page.south west) rectangle
    ([yshift=2.3cm] current page.south east);
\end{tikzpicture}

%──────────────────────────────────────────────
%  LOGOS (en-tête)
%──────────────────────────────────────────────
\vspace*{0.3cm}
\begin{minipage}[t]{0.45\textwidth}
  \hspace{0.8cm}
  \begin{minipage}[c]{0.85\textwidth}
    \placeholder{Logo École\\(figures/logo\_school.png)}
    % \includegraphics[height=2.0cm]{figures/logo_school}
  \end{minipage}
\end{minipage}
\hfill
\begin{minipage}[t]{0.45\textwidth}
  \raggedleft
  \placeholder{Logo Entreprise\\(figures/logo\_company.png)}
  % \includegraphics[height=2.0cm]{figures/logo_company}
\end{minipage}

\vspace{2cm}

%──────────────────────────────────────────────
%  TITRE PRINCIPAL
%──────────────────────────────────────────────
\begin{center}
  {\sffamily\Large\bfseries\color{ChapBlue}
    Rapport de Projet de Fin d'Année}\\[0.3cm]
  {\sffamily\normalsize\color{SecGrey}
    Stage Professionnel — \AcademicYear}

  \vspace{1.2cm}

  \begin{tcolorbox}[
    enhanced, width=13cm,
    colback=ChapBlue, colframe=AccentBlue,
    arc=6pt, fontupper=\sffamily\bfseries\large,
    halign=center, valign=center,
    top=12pt, bottom=12pt
  ]
    \textcolor{white}{\ProjectTitle}
  \end{tcolorbox}

  \vspace{0.8cm}
  {\color{SecGrey}\sffamily\normalsize
    \CompanyName\ — \CompanyCity}
\end{center}

\vspace{1.5cm}

%──────────────────────────────────────────────
%  IDENTITÉ ÉTUDIANT
%──────────────────────────────────────────────
\begin{center}
  \begin{tcolorbox}[
    enhanced, width=12cm,
    colback=LightBg, colframe=ChapBlue,
    title={\sffamily\bfseries Réalisé par},
    attach boxed title to top left={xshift=10pt,yshift=-4pt},
    boxed title style={colback=ChapBlue,colframe=ChapBlue}
  ]
    \begin{center}
      {\large\bfseries\color{ChapBlue} \StudentName}\\[4pt]
      {\color{SecGrey} \StudentField\ — \StudentYear}
    \end{center}
  \end{tcolorbox}
\end{center}

\vspace{1.2cm}

%──────────────────────────────────────────────
%  ENCADRANTS
%──────────────────────────────────────────────
\begin{center}
  \begin{tabular}{@{}ll@{}}
    \toprule
    \multicolumn{2}{c}{\sffamily\bfseries\color{ChapBlue} Encadrement} \\
    \midrule
    {\sffamily\bfseries\color{SecGrey} Encadrant professionnel :}
      & {\color{ChapBlue} \SupervisorPro} \\[4pt]
    {\sffamily\bfseries\color{SecGrey} Encadrant académique :}
      & {\color{ChapBlue} \SupervisorAca} \\
    \bottomrule
  \end{tabular}
\end{center}

\vfill

%──────────────────────────────────────────────
%  PIED DE GARDE
%──────────────────────────────────────────────
\begin{center}
  {\color{white}\sffamily\small
    \SchoolName\ \guillemotright\ Année universitaire \AcademicYear
    \guillemotright\ \SubmissionDate}
\end{center}

\end{titlepage}
\clearpage

%==========================================================
\begin{titlepage}
\thispagestyle{empty}

% Couche TikZ pour les éléments graphiques de fond
\begin{tikzpicture}[remember picture, overlay]
  % Bandeau supérieur bleu foncé
  \fill[ChapBlue]
    (current page.north west) rectangle
    ([yshift=-4.5cm] current page.north east);
  % Liseré accent
  \fill[AccentBlue]
    ([yshift=-4.5cm] current page.north west) rectangle
    ([yshift=-4.8cm] current page.north east);
  % Barre latérale gauche
  \fill[ChapBlue!15]
    (current page.south west) rectangle
    ([xshift=0.7cm] current page.north west);
  % Bandeau bas
  \fill[ChapBlue]
    (current page.south west) rectangle
    ([yshift=2cm] current page.south east);
  \fill[AccentBlue]
    ([yshift=2cm] current page.south west) rectangle
    ([yshift=2.3cm] current page.south east);
\end{tikzpicture}

%──────────────────────────────────────────────
%  LOGOS (en-tête)
%──────────────────────────────────────────────
\vspace*{0.3cm}
\begin{minipage}[t]{0.45\textwidth}
  \hspace{0.8cm}
  \begin{minipage}[c]{0.85\textwidth}
    \placeholder{Logo École\\(figures/logo\_school.png)}
    % \includegraphics[height=2.0cm]{figures/logo_school}
  \end{minipage}
\end{minipage}
\hfill
\begin{minipage}[t]{0.45\textwidth}
  \raggedleft
  \placeholder{Logo Entreprise\\(figures/logo\_company.png)}
  % \includegraphics[height=2.0cm]{figures/logo_company}
\end{minipage}

\vspace{2cm}

%──────────────────────────────────────────────
%  TITRE PRINCIPAL
%──────────────────────────────────────────────
\begin{center}
  {\sffamily\Large\bfseries\color{ChapBlue}
    Rapport de Projet de Fin d'Année}\\[0.3cm]
  {\sffamily\normalsize\color{SecGrey}
    Stage Professionnel — \AcademicYear}

  \vspace{1.2cm}

  \begin{tcolorbox}[
    enhanced, width=13cm,
    colback=ChapBlue, colframe=AccentBlue,
    arc=6pt, fontupper=\sffamily\bfseries\large,
    halign=center, valign=center,
    top=12pt, bottom=12pt
  ]
    \textcolor{white}{\ProjectTitle}
  \end{tcolorbox}

  \vspace{0.8cm}
  {\color{SecGrey}\sffamily\normalsize
    \CompanyName\ — \CompanyCity}
\end{center}

\vspace{1.5cm}

%──────────────────────────────────────────────
%  IDENTITÉ ÉTUDIANT
%──────────────────────────────────────────────
\begin{center}
  \begin{tcolorbox}[
    enhanced, width=12cm,
    colback=LightBg, colframe=ChapBlue,
    title={\sffamily\bfseries Réalisé par},
    attach boxed title to top left={xshift=10pt,yshift=-4pt},
    boxed title style={colback=ChapBlue,colframe=ChapBlue}
  ]
    \begin{center}
      {\large\bfseries\color{ChapBlue} \StudentName}\\[4pt]
      {\color{SecGrey} \StudentField\ — \StudentYear}
    \end{center}
  \end{tcolorbox}
\end{center}

\vspace{1.2cm}

%──────────────────────────────────────────────
%  ENCADRANTS
%──────────────────────────────────────────────
\begin{center}
  \begin{tabular}{@{}ll@{}}
    \toprule
    \multicolumn{2}{c}{\sffamily\bfseries\color{ChapBlue} Encadrement} \\
    \midrule
    {\sffamily\bfseries\color{SecGrey} Encadrant professionnel :}
      & {\color{ChapBlue} \SupervisorPro} \\[4pt]
    {\sffamily\bfseries\color{SecGrey} Encadrant académique :}
      & {\color{ChapBlue} \SupervisorAca} \\
    \bottomrule
  \end{tabular}
\end{center}

\vfill

%──────────────────────────────────────────────
%  PIED DE GARDE
%──────────────────────────────────────────────
\begin{center}
  {\color{white}\sffamily\small
    \SchoolName\ \guillemotright\ Année universitaire \AcademicYear
    \guillemotright\ \SubmissionDate}
\end{center}

\end{titlepage}
\clearpage

%==========================================================
\begin{titlepage}
\thispagestyle{empty}

% Couche TikZ pour les éléments graphiques de fond
\begin{tikzpicture}[remember picture, overlay]
  % Bandeau supérieur bleu foncé
  \fill[ChapBlue]
    (current page.north west) rectangle
    ([yshift=-4.5cm] current page.north east);
  % Liseré accent
  \fill[AccentBlue]
    ([yshift=-4.5cm] current page.north west) rectangle
    ([yshift=-4.8cm] current page.north east);
  % Barre latérale gauche
  \fill[ChapBlue!15]
    (current page.south west) rectangle
    ([xshift=0.7cm] current page.north west);
  % Bandeau bas
  \fill[ChapBlue]
    (current page.south west) rectangle
    ([yshift=2cm] current page.south east);
  \fill[AccentBlue]
    ([yshift=2cm] current page.south west) rectangle
    ([yshift=2.3cm] current page.south east);
\end{tikzpicture}

%──────────────────────────────────────────────
%  LOGOS (en-tête)
%──────────────────────────────────────────────
\vspace*{0.3cm}
\begin{minipage}[t]{0.45\textwidth}
  \hspace{0.8cm}
  \begin{minipage}[c]{0.85\textwidth}
    \placeholder{Logo École\\(figures/logo\_school.png)}
    % \includegraphics[height=2.0cm]{figures/logo_school}
  \end{minipage}
\end{minipage}
\hfill
\begin{minipage}[t]{0.45\textwidth}
  \raggedleft
  \placeholder{Logo Entreprise\\(figures/logo\_company.png)}
  % \includegraphics[height=2.0cm]{figures/logo_company}
\end{minipage}

\vspace{2cm}

%──────────────────────────────────────────────
%  TITRE PRINCIPAL
%──────────────────────────────────────────────
\begin{center}
  {\sffamily\Large\bfseries\color{ChapBlue}
    Rapport de Projet de Fin d'Année}\\[0.3cm]
  {\sffamily\normalsize\color{SecGrey}
    Stage Professionnel — \AcademicYear}

  \vspace{1.2cm}

  \begin{tcolorbox}[
    enhanced, width=13cm,
    colback=ChapBlue, colframe=AccentBlue,
    arc=6pt, fontupper=\sffamily\bfseries\large,
    halign=center, valign=center,
    top=12pt, bottom=12pt
  ]
    \textcolor{white}{\ProjectTitle}
  \end{tcolorbox}

  \vspace{0.8cm}
  {\color{SecGrey}\sffamily\normalsize
    \CompanyName\ — \CompanyCity}
\end{center}

\vspace{1.5cm}

%──────────────────────────────────────────────
%  IDENTITÉ ÉTUDIANT
%──────────────────────────────────────────────
\begin{center}
  \begin{tcolorbox}[
    enhanced, width=12cm,
    colback=LightBg, colframe=ChapBlue,
    title={\sffamily\bfseries Réalisé par},
    attach boxed title to top left={xshift=10pt,yshift=-4pt},
    boxed title style={colback=ChapBlue,colframe=ChapBlue}
  ]
    \begin{center}
      {\large\bfseries\color{ChapBlue} \StudentName}\\[4pt]
      {\color{SecGrey} \StudentField\ — \StudentYear}
    \end{center}
  \end{tcolorbox}
\end{center}

\vspace{1.2cm}

%──────────────────────────────────────────────
%  ENCADRANTS
%──────────────────────────────────────────────
\begin{center}
  \begin{tabular}{@{}ll@{}}
    \toprule
    \multicolumn{2}{c}{\sffamily\bfseries\color{ChapBlue} Encadrement} \\
    \midrule
    {\sffamily\bfseries\color{SecGrey} Encadrant professionnel :}
      & {\color{ChapBlue} \SupervisorPro} \\[4pt]
    {\sffamily\bfseries\color{SecGrey} Encadrant académique :}
      & {\color{ChapBlue} \SupervisorAca} \\
    \bottomrule
  \end{tabular}
\end{center}

\vfill

%──────────────────────────────────────────────
%  PIED DE GARDE
%──────────────────────────────────────────────
\begin{center}
  {\color{white}\sffamily\small
    \SchoolName\ \guillemotright\ Année universitaire \AcademicYear
    \guillemotright\ \SubmissionDate}
\end{center}

\end{titlepage}
\clearpage


%──────────────────────────────────────────────────────────
%  PAGES LIMINAIRES
%──────────────────────────────────────────────────────────
\pagenumbering{roman}
\setcounter{page}{1}

%==========================================================
%  CHAPITRE 01 — REMERCIEMENTS
%==========================================================
\chapter*{Remerciements}
\addcontentsline{toc}{chapter}{Remerciements}
\markboth{Remerciements}{Remerciements}

\begin{flushright}
  \textit{\todo{Lieu}, le \todo{date}}
\end{flushright}

\vspace{1em}

Je tiens tout d'abord à exprimer ma profonde gratitude à
\textbf{\SupervisorPro}, \todo{titre et poste}, pour son accueil
au sein de \textbf{\CompanyName}, sa disponibilité, ses précieux
conseils et la confiance qu'il/elle m'a accordée tout au long de
ce stage.

\vspace{0.8em}

Je remercie également \textbf{\SupervisorAca}, mon encadrant
académique à \textbf{\SchoolName}, pour son suivi régulier,
ses orientations méthodologiques et son soutien dans la
rédaction de ce rapport.

\vspace{0.8em}

Mes remerciements vont aussi à :
\begin{itemize}
  \item L'ensemble de l'équipe de \todo{nom du département/équipe},
        pour son accueil chaleureux et l'ambiance de travail
        enrichissante.
  \item \todo{Nom(s) de collègues ou personnes ayant aidé},
        pour leur aide technique et leurs précieux échanges.
  \item \todo{Autres personnes à remercier — direction, RH, etc.}
\end{itemize}

\vspace{0.8em}

Je souhaite également témoigner ma reconnaissance envers le
corps enseignant de \textbf{\SchoolName}, dont la formation
solide m'a permis d'aborder ce stage avec les outils nécessaires
à sa réussite.

\vspace{0.8em}

Enfin, je remercie ma famille et mes proches pour leur soutien
indéfectible et leurs encouragements tout au long de mes études.

\vspace{2em}

\begin{flushright}
  \textit{\StudentName}
\end{flushright}

\clearpage

%==========================================================
%  CHAPITRE 02 — RÉSUMÉ / ABSTRACT
%==========================================================

%──────────────────────────────────────────────────────────
%  RÉSUMÉ (Français)
%──────────────────────────────────────────────────────────
\chapter*{Résumé}
\addcontentsline{toc}{chapter}{Résumé}
\markboth{Résumé}{Résumé}

\begin{otherlanguage}{french}

\begin{tcolorbox}[
  enhanced, breakable,
  colback=LightBg, colframe=ChapBlue,
  title={\sffamily\bfseries Résumé — \ProjectTitle},
  attach boxed title to top left={xshift=8pt,yshift=-4pt},
  boxed title style={colback=ChapBlue}
]

Ce rapport présente le travail réalisé dans le cadre d'un stage de
\todo{durée, ex : 3 mois} au sein de \textbf{\CompanyName},
du \InternshipStart\ au \InternshipEnd.

\vspace{0.5em}

\todo{Rédiger ici le résumé en français en 150--250 mots. Décrire :
(1) le contexte et la problématique de l'entreprise,
(2) la solution développée et la méthodologie adoptée,
(3) les principaux résultats et livrables obtenus,
(4) les apports personnels et les perspectives.}

\vspace{0.5em}

\noindent\textbf{Mots-clés :}
\todo{mot-clé 1}, \todo{mot-clé 2}, \todo{mot-clé 3},
\todo{mot-clé 4}, \todo{mot-clé 5}.

\end{tcolorbox}
\end{otherlanguage}

\vspace{2em}

%──────────────────────────────────────────────────────────
%  ABSTRACT (English)
%──────────────────────────────────────────────────────────
\chapter*{Abstract}
\addcontentsline{toc}{chapter}{Abstract}
\markboth{Abstract}{Abstract}

\begin{otherlanguage}{english}

\begin{tcolorbox}[
  enhanced, breakable,
  colback=LightBg, colframe=AccentBlue,
  title={\sffamily\bfseries Abstract — \ProjectTitle},
  attach boxed title to top left={xshift=8pt,yshift=-4pt},
  boxed title style={colback=AccentBlue}
]

This report presents the work carried out during a \todo{duration, e.g., 3-month}
internship at \textbf{\CompanyName}, from \InternshipStart\ to \InternshipEnd.

\vspace{0.5em}

\todo{Write a 150--250 word English abstract here. Cover:
(1) the business context and problem statement,
(2) the developed solution and methodology,
(3) key results and deliverables,
(4) personal contributions and future perspectives.}

\vspace{0.5em}

\noindent\textbf{Keywords:}
\todo{keyword 1}, \todo{keyword 2}, \todo{keyword 3},
\todo{keyword 4}, \todo{keyword 5}.

\end{tcolorbox}
\end{otherlanguage}

\clearpage


%──────────────────────────────────────────────────────────
%  TABLES DES MATIÈRES / FIGURES / TABLEAUX
%──────────────────────────────────────────────────────────
{
  \hypersetup{linkcolor=black}
  \tableofcontents
  \clearpage
  \listoffigures
  \addcontentsline{toc}{chapter}{Liste des figures}
  \clearpage
  \listoftables
  \addcontentsline{toc}{chapter}{Liste des tableaux}
  \clearpage
  \lstlistoflistings
  \addcontentsline{toc}{chapter}{Liste des codes sources}
}

%──────────────────────────────────────────────────────────
%  CORPS DU RAPPORT
%──────────────────────────────────────────────────────────
\cleardoublepage
\pagenumbering{arabic}
\setcounter{page}{1}

%==========================================================
%  CHAPITRE 03 — INTRODUCTION GÉNÉRALE
%==========================================================
\chapter{Introduction générale}
\label{chap:intro}

\begin{flushright}
  \textit{\guillemotleft\ \todo{Citation inspirante en lien avec le projet}
  \guillemotright}\\
  \vspace{2pt}
  {\small\color{SecGrey}— \todo{Auteur de la citation}}
\end{flushright}

\vspace{1em}

%──────────────────────────────────────────────────────────
\section{Contexte général}
\label{sec:intro:contexte}
%──────────────────────────────────────────────────────────

\todo{Décrire en 2--3 paragraphes le contexte macro-économique,
sectoriel ou technologique dans lequel s'inscrit le projet.
Exemple : transformation numérique du secteur X, adoption croissante
du cloud, essor des solutions mobiles, etc.}

L'essor du numérique impose aux entreprises de tous secteurs une
transformation profonde de leurs systèmes d'information. Dans ce contexte,
\textbf{\CompanyName} \todo{décrire l'activité et les enjeux spécifiques}.

\vspace{0.5em}

\todo{Paragraphe 2 : préciser le périmètre du stage et l'équipe d'accueil.}

%──────────────────────────────────────────────────────────
\section{Problématique}
\label{sec:intro:problematique}
%──────────────────────────────────────────────────────────

\todo{Formuler clairement la problématique à laquelle répond le projet.
Une problématique est une question précise, pas un constat vague.
Exemple : « Comment automatiser le processus de validation des commandes
afin de réduire le délai de traitement de 48 h à moins de 2 h, tout en
garantissant la traçabilité des opérations ? »}

Face à ces enjeux, la question centrale à laquelle ce projet tente de
répondre est la suivante :

\begin{important}
  \centering
  \textit{« \todo{Formuler la problématique principale} »}
\end{important}

%──────────────────────────────────────────────────────────
\section{Objectifs du projet}
\label{sec:intro:objectifs}
%──────────────────────────────────────────────────────────

Ce projet de fin d'année vise à atteindre les objectifs suivants :

\begin{enumerate}
  \item \textbf{Objectif fonctionnel :}
        \todo{ex : développer une interface permettant la saisie et le
        suivi des commandes en temps réel.}

  \item \textbf{Objectif technique :}
        \todo{ex : concevoir une architecture REST scalable avec
        \tech{Spring Boot} et \tech{React}.}

  \item \textbf{Objectif qualité :}
        \todo{ex : atteindre une couverture de tests supérieure à 75\%
        et respecter les standards de codage de l'équipe.}

  \item \textbf{Objectif de déploiement :}
        \todo{ex : mettre en production sur l'infrastructure Docker
        de l'entreprise avant la fin du stage.}
\end{enumerate}

%──────────────────────────────────────────────────────────
\section{Méthodologie de travail}
\label{sec:intro:methodologie}
%──────────────────────────────────────────────────────────

\todo{Présenter brièvement la méthodologie adoptée (Agile/Scrum, cycle
en V, Kanban, etc.) et justifier ce choix au regard des contraintes du projet.}

Le projet a été conduit selon la méthodologie \textbf{\todo{Scrum / Agile / Cycle en V}},
en raison de \todo{justification : besoins évolutifs, petite équipe, délais courts, etc.}.
Les réunions hebdomadaires avec l'encadrant professionnel et les revues
de sprint bimensuelles ont permis d'assurer un suivi rigoureux.

%──────────────────────────────────────────────────────────
\section{Structure du rapport}
\label{sec:intro:structure}
%──────────────────────────────────────────────────────────

Ce rapport est organisé en \textbf{cinq chapitres principaux} :

\begin{description}[leftmargin=!, labelwidth=\widthof{\textbf{Chapitre 5}}]
  \item[\textbf{Chapitre 1}] \textit{Présentation de l'organisme d'accueil} —
    présente \CompanyName, son activité, son organigramme et
    l'environnement de travail.

  \item[\textbf{Chapitre 2}] \textit{Contexte général du projet} —
    analyse l'existant, détaille les besoins fonctionnels et
    non-fonctionnels et précise le cahier des charges.

  \item[\textbf{Chapitre 3}] \textit{Analyse et conception} —
    expose la méthodologie, les diagrammes UML et l'architecture
    technique retenue.

  \item[\textbf{Chapitre 4}] \textit{Réalisation} —
    détaille l'implémentation, les technologies utilisées et
    présente les principales interfaces.

  \item[\textbf{Chapitre 5}] \textit{Tests et résultats} —
    décrit la stratégie de tests, les résultats obtenus et
    les performances mesurées.
\end{description}

Une conclusion générale clôt le rapport en dressant le bilan
du projet et en proposant des perspectives d'évolution.

\clearpage

%==========================================================
%  CHAPITRE 1 — PRÉSENTATION DE L'ORGANISME D'ACCUEIL
%==========================================================
\chapter{Présentation de l'organisme d'accueil}
\label{chap:organisme}

\begin{remarque}
  Ce chapitre dresse un portrait de \textbf{\CompanyName},
  en présentant ses activités, son positionnement, son
  organisation interne et le département au sein duquel
  le stage s'est déroulé.
\end{remarque}

%──────────────────────────────────────────────────────────
\section{Présentation générale de l'entreprise}
\label{sec:org:presentation}
%──────────────────────────────────────────────────────────

\subsection{Fiche d'identité}

\begin{center}
\begin{tabular}{>{\bfseries\sffamily\color{SecGrey}}l l}
  \toprule
  Raison sociale       & \CompanyName \\
  Secteur d'activité   & \todo{ex : Informatique / Finance / Industrie} \\
  Forme juridique      & \todo{ex : SA, SARL, SAS, Établissement public} \\
  Date de création     & \todo{ex : 1995} \\
  Siège social         & \todo{Adresse complète} \\
  Site Web             & \url{\todo{https://www.entreprise.com}} \\
  Effectif             & \todo{ex : 250 employés} \\
  Chiffre d'affaires   & \todo{ex : XX M€ (20XX)} \\
  \bottomrule
\end{tabular}
\end{center}

\subsection{Activités et domaines d'expertise}

\todo{Décrire en 2--3 paragraphes les activités principales de
l'entreprise, ses clients, ses marchés et sa vision stratégique.
Exemple : « \CompanyName\ est spécialisée dans \ldots\ Elle propose
une gamme de solutions couvrant \ldots\ et sert une clientèle de
\ldots »}

\subsection{Présence géographique et chiffres clés}

\todo{Mentionner les implantations (nationales ou internationales),
les certifications (ISO, etc.), les partenariats ou labels obtenus.}

%──────────────────────────────────────────────────────────
\section{Organigramme et organisation interne}
\label{sec:org:organigramme}
%──────────────────────────────────────────────────────────

La figure~\ref{fig:organigramme} présente l'organigramme général
de \textbf{\CompanyName}.

\begin{figure}[H]
  \centering
  %─── Organigramme TikZ ────────────────────────────────
  \begin{tikzpicture}[
    node distance=1.5cm and 2cm,
    box/.style={
      rectangle, rounded corners=4pt,
      fill=ChapBlue!12, draw=ChapBlue, text=ChapBlue,
      font=\sffamily\small\bfseries,
      minimum width=3.2cm, minimum height=0.9cm,
      align=center, drop shadow
    },
    line/.style={draw=ChapBlue!60, -{Stealth[length=6pt]}, thick}
  ]
    % Direction
    \node[box, fill=ChapBlue, text=white] (DG) {Direction Générale};

    % Niveau 2
    \node[box, below left=of DG, xshift=-1cm]  (DIR1) {\todo{Direction 1}};
    \node[box, below=of DG]                    (DIR2) {\todo{Direction 2}};
    \node[box, below right=of DG, xshift=1cm]  (DIR3) {\todo{Direction 3}};

    % Niveau 3 (exemple sous DIR2)
    \node[box, below left=of DIR2]   (D21) {\todo{Dépt. A}};
    \node[box, below right=of DIR2]  (D22) {\todo{Dépt. B}};

    % Flèches
    \draw[line] (DG) -- (DIR1);
    \draw[line] (DG) -- (DIR2);
    \draw[line] (DG) -- (DIR3);
    \draw[line] (DIR2) -- (D21);
    \draw[line] (DIR2) -- (D22);
  \end{tikzpicture}
  \caption{Organigramme simplifié de \CompanyName\ \todo{(à adapter)}}
  \label{fig:organigramme}
\end{figure}

\todo{Commenter l'organigramme en 3--5 phrases : nombre de directions,
hiérarchie, place du département d'accueil.}

%──────────────────────────────────────────────────────────
\section{Département d'accueil et équipe}
\label{sec:org:departement}
%──────────────────────────────────────────────────────────

Le stage s'est déroulé au sein du département
\textbf{\todo{Nom du département}}, rattaché à la direction
\textbf{\todo{Nom de la direction}}.

\subsection{Missions du département}

\begin{itemize}
  \item \todo{Mission principale du département}
  \item \todo{Mission secondaire 1}
  \item \todo{Mission secondaire 2}
  \item \todo{Outils et systèmes gérés par le département}
\end{itemize}

\subsection{Composition de l'équipe}

\begin{center}
\begin{tabular}{>{\sffamily\bfseries}l l l}
  \toprule
  \textbf{Nom / Prénom} & \textbf{Rôle} & \textbf{Lien avec le stage} \\
  \midrule
  \SupervisorPro         & \todo{Titre/Poste}    & Encadrant direct \\
  \todo{Collègue 1}      & \todo{Développeur}    & Collaboration technique \\
  \todo{Collègue 2}      & \todo{Analyste}       & Revue fonctionnelle \\
  \todo{Collègue 3}      & \todo{DevOps}         & Déploiement \\
  \bottomrule
\end{tabular}
\end{center}

%──────────────────────────────────────────────────────────
\section{Environnement de travail}
\label{sec:org:environnement}
%──────────────────────────────────────────────────────────

\todo{Décrire brièvement les conditions de travail : locaux, horaires,
outils mis à disposition (matériel, licences, VPN, accès serveurs),
culture d'entreprise.}

\begin{itemize}
  \item \textbf{Infrastructure :} \todo{ex : Serveurs on-premise + Cloud AWS}
  \item \textbf{Outils collaboratifs :} \todo{ex : Slack, Teams, Confluence}
  \item \textbf{Environnement de développement :} \todo{ex : Linux Ubuntu 22.04 LTS}
  \item \textbf{Méthodologie de travail :} \todo{ex : Agile Scrum, sprints de 2 semaines}
\end{itemize}

\clearpage

%==========================================================
%  CHAPITRE 2 — CONTEXTE GÉNÉRAL DU PROJET
%==========================================================
\chapter{Contexte général du projet}
\label{chap:contexte}

%──────────────────────────────────────────────────────────
\section{Étude de l'existant}
\label{sec:ctx:existant}
%──────────────────────────────────────────────────────────

\subsection{Situation actuelle}

\todo{Décrire le système ou le processus existant avant le projet :
outils utilisés, fonctionnement manuel ou semi-automatisé, acteurs
impliqués, flux d'information.}

\todo{Si applicable, insérer un schéma du processus existant (diagramme BPMN,
flowchart) avec \textbackslash placeholder\{...\}.}

\placeholder{Schéma du processus existant (ex : diagramme BPMN, flowchart —
à insérer dans figures/existant.png)}

\subsection{Limites et problèmes de l'existant}

\begin{center}
\begin{tabular}{>{\sffamily\bfseries}p{4cm} p{9cm}}
  \toprule
  \textbf{Problème identifié} & \textbf{Impact} \\
  \midrule
  \todo{Saisie manuelle}
    & \todo{Risque d'erreurs, perte de temps (ex : 2h/jour)} \\
  \todo{Absence de traçabilité}
    & \todo{Impossible de retrouver l'historique des opérations} \\
  \todo{Silos de données}
    & \todo{Données dispersées dans plusieurs fichiers Excel non synchronisés} \\
  \todo{Scalabilité limitée}
    & \todo{Le système actuel ne supporte pas une charge > N utilisateurs simultanés} \\
  \bottomrule
\end{tabular}
\end{center}

%──────────────────────────────────────────────────────────
\section{Analyse des besoins}
\label{sec:ctx:besoins}
%──────────────────────────────────────────────────────────

\subsection{Identification des acteurs}

\begin{center}
\begin{tabular}{>{\sffamily\bfseries}l l p{7cm}}
  \toprule
  \textbf{Acteur} & \textbf{Type} & \textbf{Rôle dans le système} \\
  \midrule
  \todo{Administrateur} & Interne & \todo{Gestion des utilisateurs et des droits} \\
  \todo{Gestionnaire}   & Interne & \todo{Saisie et validation des données métier} \\
  \todo{Utilisateur}    & Interne & \todo{Consultation des tableaux de bord} \\
  \todo{Client / API tierce} & Externe & \todo{Déclenchement d'actions via API} \\
  \bottomrule
\end{tabular}
\end{center}

\subsection{Besoins fonctionnels}

\begin{enumerate}
  \item \textbf{BF-01 :} \todo{Description du besoin fonctionnel 1.
        Exemple : Le système doit permettre à un gestionnaire de créer,
        modifier et supprimer une commande.}
  \item \textbf{BF-02 :} \todo{Description du besoin fonctionnel 2.}
  \item \textbf{BF-03 :} \todo{Description du besoin fonctionnel 3.}
  \item \textbf{BF-04 :} \todo{Description du besoin fonctionnel 4.}
  \item \textbf{BF-05 :} \todo{Description du besoin fonctionnel 5.}
  \item \todo{Ajouter d'autres besoins fonctionnels.}
\end{enumerate}

\subsection{Besoins non fonctionnels}

\begin{itemize}
  \item \textbf{Performance :} \todo{ex : temps de réponse < 300 ms pour 95\% des requêtes}
  \item \textbf{Disponibilité :} \todo{ex : SLA 99,5\% (hors maintenance planifiée)}
  \item \textbf{Sécurité :} \todo{ex : authentification JWT, chiffrement TLS 1.3, RBAC}
  \item \textbf{Maintenabilité :} \todo{ex : couverture de tests > 75\%, documentation Swagger}
  \item \textbf{Scalabilité :} \todo{ex : architecture permettant la montée en charge horizontale}
  \item \textbf{Ergonomie :} \todo{ex : interface responsive, accessible WCAG 2.1 niveau AA}
\end{itemize}

%──────────────────────────────────────────────────────────
\section{Cahier des charges}
\label{sec:ctx:cdc}
%──────────────────────────────────────────────────────────

\subsection{Périmètre du projet}

\todo{Délimiter précisément ce qui est dans le périmètre (in-scope)
et hors périmètre (out-of-scope) du projet.}

\begin{center}
\begin{tabularx}{\textwidth}{X X}
  \toprule
  \rowcolor{ChapBlue!10}
  \textbf{Dans le périmètre (In-scope)} & \textbf{Hors périmètre (Out-of-scope)} \\
  \midrule
  \todo{Fonctionnalité A} & \todo{Module existant non touché} \\
  \todo{Fonctionnalité B} & \todo{Intégration ERP tiers} \\
  \todo{Fonctionnalité C} & \todo{Application mobile} \\
  \bottomrule
\end{tabularx}
\end{center}

\subsection{Contraintes du projet}

\begin{itemize}
  \item \textbf{Contraintes techniques :} \todo{stack imposée, compatibilité
        avec l'infrastructure existante, langage imposé par l'équipe}
  \item \textbf{Contraintes temporelles :} \todo{livraison V1 attendue à J+6 semaines}
  \item \textbf{Contraintes budgétaires :} \todo{utilisation de solutions open-source uniquement}
  \item \textbf{Contraintes réglementaires :} \todo{conformité RGPD, ISO 27001, etc.}
\end{itemize}

%──────────────────────────────────────────────────────────
\section{Solutions étudiées et choix technologiques}
\label{sec:ctx:solutions}
%──────────────────────────────────────────────────────────

\todo{Si plusieurs solutions ont été étudiées avant de choisir la
solution finale, présenter ici le benchmark ou la matrice de décision.}

\begin{center}
\begin{tabular}{>{\sffamily\bfseries}l c c c c c}
  \toprule
  \rowcolor{ChapBlue!10}
  \textbf{Critère / Solution} & \textbf{Poids}
    & \textbf{\todo{Sol. A}} & \textbf{\todo{Sol. B}} & \textbf{\todo{Sol. C}}
    & \textbf{Retenu} \\
  \midrule
  Performance   & 30\% & 4/5 & 3/5 & 5/5 & \todo{Sol. C} \\
  Coût          & 20\% & 5/5 & 3/5 & 4/5 & \\
  Maintenabilité& 25\% & 3/5 & 5/5 & 4/5 & \\
  Communauté    & 15\% & 4/5 & 4/5 & 5/5 & \\
  Intégration   & 10\% & 5/5 & 2/5 & 4/5 & \\
  \midrule
  \textbf{Score pondéré} & & \todo{XX} & \todo{XX} & \todo{XX} & \\
  \bottomrule
\end{tabular}
\end{center}

\clearpage

%==========================================================
%  CHAPITRE 3 — ANALYSE ET CONCEPTION
%==========================================================
\chapter{Analyse et conception}
\label{chap:conception}

%──────────────────────────────────────────────────────────
\section{Méthodologie de développement}
\label{sec:cpt:methodologie}
%──────────────────────────────────────────────────────────

\todo{Décrire la méthodologie adoptée (Scrum, Kanban, Cycle en V,
RUP, etc.) avec ses artefacts, ses rôles et ses rituels.}

\subsection{Processus Scrum}

\begin{center}
% Schéma simplifié du cycle Scrum
\begin{tikzpicture}[
  node distance=1.5cm,
  box/.style={
    rectangle, rounded corners=5pt,
    fill=ChapBlue!12, draw=ChapBlue,
    font=\sffamily\footnotesize\bfseries, text=ChapBlue,
    minimum width=2.8cm, minimum height=0.85cm,
    align=center
  },
  arrow/.style={-{Stealth[length=5pt]}, thick, draw=AccentBlue}
]
  \node[box, fill=ChapBlue, text=white] (PB)  {Product\\Backlog};
  \node[box, right=of PB]               (SP)  {Sprint\\Planning};
  \node[box, right=of SP]               (SB)  {Sprint\\Backlog};
  \node[box, right=of SB]               (DEV) {Sprint\\(2 sem.)};
  \node[box, right=of DEV]              (REV) {Sprint\\Review};
  \node[box, below=of DEV]              (RET) {Rétrospective};

  \draw[arrow] (PB) -- (SP);
  \draw[arrow] (SP) -- (SB);
  \draw[arrow] (SB) -- (DEV);
  \draw[arrow] (DEV) -- (REV);
  \draw[arrow] (REV) -- ++(0,-1.2) -| (SP);
  \draw[arrow] (DEV) -- (RET);
\end{tikzpicture}
\end{center}

\vspace{0.5em}

\begin{center}
\begin{tabular}{>{\sffamily\bfseries}l l p{7cm}}
  \toprule
  \textbf{Artéfact / Rôle} & \textbf{Responsable} & \textbf{Description} \\
  \midrule
  Product Backlog & \todo{Product Owner} & Liste priorisée de toutes les user stories \\
  Sprint Backlog  & Équipe              & Sélection de stories à réaliser dans le sprint \\
  Daily Standup   & Équipe              & Réunion quotidienne de 15 min \\
  Sprint Review   & Équipe + Stakeholders & Démo de l'incrément livré \\
  Rétrospective   & Équipe + Scrum Master & Amélioration continue du processus \\
  \bottomrule
\end{tabular}
\end{center}

%──────────────────────────────────────────────────────────
\section{Diagrammes UML}
\label{sec:cpt:uml}
%──────────────────────────────────────────────────────────

\subsection{Diagramme de cas d'utilisation}

\placeholder{Diagramme de cas d'utilisation global\\
(à insérer dans figures/usecase.png ou à générer avec PlantUML)}

La figure~\ref{fig:usecase} présente les interactions entre les
différents acteurs et le système.

\begin{figure}[H]
  \centering
  \placeholder{Diagramme de cas d'utilisation — figures/usecase.png}
  % \includegraphics[width=0.85\textwidth]{figures/usecase}
  \caption{Diagramme de cas d'utilisation global}
  \label{fig:usecase}
\end{figure}

\todo{Commenter chaque groupe de cas d'utilisation en précisant
l'acteur concerné et le déclencheur de l'interaction.}

\subsection{Diagramme de classes}

\begin{figure}[H]
  \centering
  \placeholder{Diagramme de classes — figures/class\_diagram.png}
  % \includegraphics[width=0.9\textwidth]{figures/class_diagram}
  \caption{Diagramme de classes du domaine}
  \label{fig:classes}
\end{figure}

\todo{Décrire les principales entités (classes), leurs attributs,
leurs méthodes et leurs relations (association, agrégation,
composition, héritage).}

\subsection{Diagramme de séquence — Scénario principal}

\begin{figure}[H]
  \centering
  \placeholder{Diagramme de séquence — scénario principal\\
               (figures/sequence\_main.png)}
  % \includegraphics[width=0.9\textwidth]{figures/sequence_main}
  \caption{Diagramme de séquence : \todo{Nom du scénario principal}}
  \label{fig:sequence}
\end{figure}

\todo{Expliquer le déroulement du scénario illustré.}

\subsection{Diagramme d'activité}

\begin{figure}[H]
  \centering
  \placeholder{Diagramme d'activité — figures/activity.png}
  % \includegraphics[width=0.75\textwidth]{figures/activity}
  \caption{Diagramme d'activité : \todo{Nom du processus métier}}
  \label{fig:activity}
\end{figure}

%──────────────────────────────────────────────────────────
\section{Architecture technique}
\label{sec:cpt:architecture}
%──────────────────────────────────────────────────────────

\subsection{Architecture globale}

\todo{Décrire l'architecture retenue (monolithique, microservices,
MVC, hexagonale, événementielle…) et justifier ce choix.}

\begin{figure}[H]
  \centering
  \begin{tikzpicture}[
    node distance=1.4cm and 2.2cm,
    layer/.style={
      rectangle, rounded corners=6pt,
      draw=#1, fill=#1!10,
      font=\sffamily\small\bfseries, text=#1,
      minimum width=4cm, minimum height=1cm,
      align=center
    },
    arrow/.style={-{Stealth[length=5pt]}, thick, draw=SecGrey}
  ]
    % Couches
    \node[layer=AccentBlue]               (frontend)  {Frontend\\(\todo{React / Angular})};
    \node[layer=ChapBlue, below=of frontend]  (backend)
                                                 {Backend API\\(\todo{Spring Boot / FastAPI})};
    \node[layer=SecGrey,  below=of backend]   (db)
                                                 {Base de données\\(\todo{PostgreSQL / MongoDB})};
    % Côté
    \node[layer=WarningOrange, right=3cm of backend] (ext)
                                                 {Services externes\\(\todo{Email, OAuth, S3})};
    % Flèches
    \draw[arrow] (frontend) -- node[right,font=\tiny\sffamily]{REST/HTTP} (backend);
    \draw[arrow] (backend)  -- node[right,font=\tiny\sffamily]{JDBC/ORM}  (db);
    \draw[arrow] (backend)  -- node[above,font=\tiny\sffamily]{API calls} (ext);
  \end{tikzpicture}
  \caption{Architecture en couches du système}
  \label{fig:architecture}
\end{figure}

\subsection{Architecture de base de données}

\todo{Présenter le schéma de la base de données (MCD / MLD) ou décrire
les collections si NoSQL.}

\begin{figure}[H]
  \centering
  \placeholder{Schéma entité-relation / MLD — figures/db\_schema.png}
  % \includegraphics[width=0.9\textwidth]{figures/db_schema}
  \caption{Schéma de la base de données}
  \label{fig:db}
\end{figure}

\subsection{Description des composants principaux}

\begin{description}[leftmargin=0pt]
  \item[\textcolor{ChapBlue}{\sffamily\bfseries Module A — \todo{Nom}}]
    \todo{Décrire le rôle, les responsabilités et les interfaces
    exposées par ce module.}

  \item[\textcolor{ChapBlue}{\sffamily\bfseries Module B — \todo{Nom}}]
    \todo{Décrire le rôle, les responsabilités et les interfaces
    exposées par ce module.}

  \item[\textcolor{ChapBlue}{\sffamily\bfseries Module C — \todo{Nom}}]
    \todo{Décrire le rôle, les responsabilités et les interfaces
    exposées par ce module.}
\end{description}

%──────────────────────────────────────────────────────────
\section{Modélisation des données}
\label{sec:cpt:donnees}
%──────────────────────────────────────────────────────────

\todo{Si applicable, décrire les principales entités du domaine,
leur modélisation en base de données et les règles de gestion associées.}

\begin{center}
\begin{tabular}{>{\sffamily\bfseries}l >{\ttfamily}l l l}
  \toprule
  \textbf{Entité} & \textbf{Table} & \textbf{Clé primaire} & \textbf{Description} \\
  \midrule
  \todo{Utilisateur} & \todo{users}    & id (UUID)  & \todo{Description} \\
  \todo{Commande}    & \todo{orders}   & id (SERIAL)& \todo{Description} \\
  \todo{Produit}     & \todo{products} & id (SERIAL)& \todo{Description} \\
  \todo{Entité 4}    & \todo{table\_4} & id         & \todo{Description} \\
  \bottomrule
\end{tabular}
\end{center}

\clearpage

%==========================================================
%  CHAPITRE 4 — RÉALISATION
%==========================================================
\chapter{Réalisation}
\label{chap:realisation}

%──────────────────────────────────────────────────────────
\section{Environnement de développement}
\label{sec:real:env}
%──────────────────────────────────────────────────────────

\begin{center}
\begin{tabular}{>{\sffamily\bfseries\color{SecGrey}}l l l}
  \toprule
  \textbf{Catégorie} & \textbf{Outil / Technologie} & \textbf{Version} \\
  \midrule
  OS de développement   & \todo{Ubuntu 22.04 / Windows 11} & — \\
  Langage back-end      & \todo{Java / Python / Node.js}   & \todo{17 / 3.11 / 18} \\
  Framework back-end    & \todo{Spring Boot / FastAPI}      & \todo{3.2 / 0.110} \\
  Langage front-end     & \todo{JavaScript / TypeScript}    & \todo{ES2022} \\
  Framework front-end   & \todo{React / Vue / Angular}      & \todo{18.2} \\
  Base de données       & \todo{PostgreSQL / MySQL}         & \todo{15.4} \\
  ORM                   & \todo{Hibernate / SQLAlchemy}     & \todo{6.x} \\
  Conteneurisation      & \todo{Docker + Compose}           & \todo{24.0} \\
  Gestion de version    & \todo{Git + GitHub / GitLab}      & — \\
  CI/CD                 & \todo{GitHub Actions / Jenkins}   & — \\
  IDE                   & \todo{IntelliJ IDEA / VS Code}    & \todo{2024.1} \\
  \bottomrule
\end{tabular}
\end{center}

%──────────────────────────────────────────────────────────
\section{Implémentation — Module \todo{A}}
\label{sec:real:moduleA}
%──────────────────────────────────────────────────────────

\subsection{Description et fonctionnalités}

\todo{Décrire les fonctionnalités développées dans ce module,
les fichiers / packages concernés et les choix d'implémentation.}

\subsection{Extrait de code significatif}

L'extrait suivant illustre \todo{ce que fait ce code — ex : la logique
de création d'une commande avec validation métier}.

\begin{lstlisting}[
  language=Java,
  caption={\todo{Titre du snippet — ex : Service de création de commande}},
  label={lst:create-order}
]
// [A COMPLETER — coller ici un extrait de code representatif]
@Service
@Transactional
public class OrderService {

    private final OrderRepository orderRepository;
    private final InventoryClient inventoryClient;

    public OrderService(OrderRepository repo, InventoryClient inv) {
        this.orderRepository = repo;
        this.inventoryClient = inv;
    }

    public OrderDTO createOrder(CreateOrderRequest request) {
        // Valider la disponibilite en stock
        boolean inStock = inventoryClient.checkStock(
            request.getProductId(), request.getQuantity()
        );
        if (!inStock) {
            throw new InsufficientStockException(
                "Produit " + request.getProductId() + " indisponible."
            );
        }
        // Persister la commande
        Order order = Order.builder()
            .product(request.getProductId())
            .quantity(request.getQuantity())
            .status(OrderStatus.PENDING)
            .createdAt(Instant.now())
            .build();
        return OrderMapper.toDTO(orderRepository.save(order));
    }
}
\end{lstlisting}

\todo{Expliquer les choix de conception : pattern Repository, injection
de dépendances, gestion des exceptions, etc.}

%──────────────────────────────────────────────────────────
\section{Implémentation — Module \todo{B}}
\label{sec:real:moduleB}
%──────────────────────────────────────────────────────────

\subsection{Description et fonctionnalités}

\todo{Décrire les fonctionnalités du module B.}

\subsection{Extrait de code — \todo{exemple Python/JavaScript}}

\begin{lstlisting}[
  language=Python,
  caption={\todo{Titre du snippet Python}},
  label={lst:python-snippet}
]
# [A COMPLETER]
from fastapi import APIRouter, HTTPException, Depends
from sqlalchemy.orm import Session
from . import schemas, crud, database

router = APIRouter(prefix="/api/v1/items", tags=["items"])

@router.get("/{item_id}", response_model=schemas.ItemResponse)
def get_item(
    item_id: int,
    db: Session = Depends(database.get_db)
):
    """Recupere un item par son identifiant."""
    db_item = crud.get_item(db, item_id=item_id)
    if db_item is None:
        raise HTTPException(status_code=404, detail="Item non trouve")
    return db_item
\end{lstlisting}

%──────────────────────────────────────────────────────────
\section{Interfaces utilisateur}
\label{sec:real:ui}
%──────────────────────────────────────────────────────────

Cette section présente les principales interfaces de l'application.

\subsection{Page d'authentification}

\begin{figure}[H]
  \centering
  \placeholder{Capture d'écran : Page de connexion — figures/ui\_login.png}
  % \includegraphics[width=0.7\textwidth]{figures/ui_login}
  \caption{Interface d'authentification}
  \label{fig:ui-login}
\end{figure}

\todo{Décrire les éléments de l'interface : champs, validation, gestion
des erreurs, redirection après authentification.}

\subsection{Tableau de bord principal}

\begin{figure}[H]
  \centering
  \placeholder{Capture d'écran : Dashboard — figures/ui\_dashboard.png}
  % \includegraphics[width=\textwidth]{figures/ui_dashboard}
  \caption{Tableau de bord principal}
  \label{fig:ui-dashboard}
\end{figure}

\todo{Décrire les KPIs affichés, les graphiques, les actions rapides
disponibles depuis le dashboard.}

\subsection{\todo{Page / Module 3}}

\begin{figure}[H]
  \centering
  \placeholder{Capture d'écran : \todo{Nom de la page} — figures/ui\_page3.png}
  % \includegraphics[width=\textwidth]{figures/ui_page3}
  \caption{\todo{Titre de la figure}}
  \label{fig:ui-page3}
\end{figure}

%──────────────────────────────────────────────────────────
\section{Configuration et déploiement}
\label{sec:real:deploiement}
%──────────────────────────────────────────────────────────

\subsection{Docker Compose}

\begin{lstlisting}[
  language={},
  caption={Extrait du fichier docker-compose.yml},
  label={lst:docker-compose}
]
# [A COMPLETER — adapter a votre projet]
version: "3.9"
services:
  backend:
    build: ./backend
    ports:
      - "8080:8080"
    environment:
      - SPRING_DATASOURCE_URL=jdbc:postgresql://db:5432/mydb
      - SPRING_DATASOURCE_USERNAME=${DB_USER}
      - SPRING_DATASOURCE_PASSWORD=${DB_PASS}
    depends_on:
      - db

  frontend:
    build: ./frontend
    ports:
      - "3000:80"
    depends_on:
      - backend

  db:
    image: postgres:15-alpine
    volumes:
      - pgdata:/var/lib/postgresql/data
    environment:
      POSTGRES_DB: mydb
      POSTGRES_USER: ${DB_USER}
      POSTGRES_PASSWORD: ${DB_PASS}

volumes:
  pgdata:
\end{lstlisting}

\subsection{Pipeline CI/CD}

\begin{figure}[H]
  \centering
  \begin{tikzpicture}[
    node distance=0.5cm and 1.2cm,
    step/.style={
      rectangle, rounded corners=4pt,
      fill=ChapBlue!12, draw=ChapBlue,
      font=\sffamily\scriptsize\bfseries, text=ChapBlue,
      minimum width=2.2cm, minimum height=0.8cm,
      align=center
    },
    arrow/.style={-{Stealth[length=4pt]}, draw=AccentBlue, thick}
  ]
    \node[step] (push)  {Git Push};
    \node[step, right=of push] (build)  {Build};
    \node[step, right=of build] (test)  {Tests};
    \node[step, right=of test]  (scan)  {Analyse\\statique};
    \node[step, right=of scan]  (image) {Docker\\Image};
    \node[step, right=of image] (deploy){Déploiement};

    \draw[arrow] (push)  -- (build);
    \draw[arrow] (build) -- (test);
    \draw[arrow] (test)  -- (scan);
    \draw[arrow] (scan)  -- (image);
    \draw[arrow] (image) -- (deploy);
  \end{tikzpicture}
  \caption{Pipeline CI/CD — \todo{GitHub Actions / Jenkins}}
  \label{fig:cicd}
\end{figure}

\todo{Décrire chaque étape du pipeline : outil utilisé, durée moyenne,
critères de succès/échec.}

\clearpage

%==========================================================
%  CHAPITRE 5 — TESTS ET RÉSULTATS
%==========================================================
\chapter{Tests et résultats}
\label{chap:tests}

%──────────────────────────────────────────────────────────
\section{Stratégie de tests}
\label{sec:tests:strategie}
%──────────────────────────────────────────────────────────

\todo{Décrire la stratégie de tests adoptée : quels types de tests,
à quel niveau, par qui et avec quels outils.}

La pyramide de tests suivante illustre la répartition adoptée :

\begin{center}
\begin{tikzpicture}
  % Pyramide de tests
  \fill[ChapBlue!80]   (0,0)   -- (6,0)   -- (3,1.5) -- cycle;
  \fill[AccentBlue!60] (0.8,1.5)-- (5.2,1.5)-- (3,3)   -- cycle;
  \fill[WarningOrange!50](1.8,3) -- (4.2,3)  -- (3,4.5) -- cycle;

  % Labels
  \node[font=\sffamily\small\bfseries, text=white] at (3,0.7)
    {Tests Unitaires (\todo{ex: JUnit 5, pytest})};
  \node[font=\sffamily\small\bfseries, text=white] at (3,2.25)
    {Tests d'Intégration (\todo{Testcontainers, REST-assured})};
  \node[font=\sffamily\footnotesize\bfseries, text=white] at (3,3.75)
    {Tests E2E\\(\todo{Cypress / Selenium})};

  % Annotations côté droit
  \node[right, font=\sffamily\footnotesize, text=SecGrey] at (6.2,0.7)
    {\todo{~120 tests}};
  \node[right, font=\sffamily\footnotesize, text=SecGrey] at (6.2,2.25)
    {\todo{~30 tests}};
  \node[right, font=\sffamily\footnotesize, text=SecGrey] at (6.2,3.75)
    {\todo{~10 scénarios}};
\end{tikzpicture}
\end{center}

%──────────────────────────────────────────────────────────
\section{Tests unitaires}
\label{sec:tests:unitaires}
%──────────────────────────────────────────────────────────

\subsection{Outillage}

\begin{itemize}
  \item \tech{\todo{JUnit 5}} — Framework de tests unitaires
  \item \tech{\todo{Mockito}} — Mock des dépendances
  \item \tech{\todo{JaCoCo}} — Mesure de la couverture de code
\end{itemize}

\subsection{Exemple de test unitaire}

\begin{lstlisting}[
  language=Java,
  caption={Test unitaire : \todo{Cas de test — ex : OrderServiceTest}},
  label={lst:unit-test}
]
// [A COMPLETER]
@ExtendWith(MockitoExtension.class)
class OrderServiceTest {

    @Mock private OrderRepository orderRepository;
    @Mock private InventoryClient inventoryClient;

    @InjectMocks private OrderService orderService;

    @Test
    @DisplayName("Devrait lever une exception si le stock est insuffisant")
    void createOrder_shouldThrow_whenStockInsufficient() {
        // Arrange
        CreateOrderRequest req = new CreateOrderRequest("PROD-01", 100);
        when(inventoryClient.checkStock("PROD-01", 100)).thenReturn(false);

        // Act & Assert
        assertThrows(
            InsufficientStockException.class,
            () -> orderService.createOrder(req)
        );
        verify(orderRepository, never()).save(any());
    }
}
\end{lstlisting}

\subsection{Résultats de couverture}

\begin{center}
\begin{tabular}{>{\sffamily\bfseries}l r r r}
  \toprule
  \textbf{Module / Package}   & \textbf{Lignes} & \textbf{Couvertes} & \textbf{Taux} \\
  \midrule
  \todo{service}              & \todo{320}  & \todo{275}  & \todo{85,9\%} \\
  \todo{repository}           & \todo{80}   & \todo{65}   & \todo{81,3\%} \\
  \todo{controller}           & \todo{150}  & \todo{110}  & \todo{73,3\%} \\
  \todo{model}                & \todo{200}  & \todo{190}  & \todo{95,0\%} \\
  \midrule
  \textbf{Total}              & \todo{750}  & \todo{640}  & \textbf{\todo{85,3\%}} \\
  \bottomrule
\end{tabular}
\end{center}

%──────────────────────────────────────────────────────────
\section{Tests d'intégration}
\label{sec:tests:integration}
%──────────────────────────────────────────────────────────

\todo{Décrire les tests d'intégration : périmètre testé (API REST,
base de données, services externes), outils (Testcontainers, WireMock,
Postman/Newman) et principaux résultats.}

\begin{center}
\begin{tabular}{>{\sffamily\bfseries}l c c c}
  \toprule
  \textbf{Scénario}              & \textbf{Méthode HTTP} & \textbf{Code attendu} & \textbf{Résultat} \\
  \midrule
  Créer une commande valide      & POST   & 201 Created    & \textcolor{green!60!black}{\bfseries PASS} \\
  Commande avec stock insuffisant& POST   & 422 Unprocessable & \textcolor{green!60!black}{\bfseries PASS} \\
  Récupérer une commande existante & GET  & 200 OK         & \textcolor{green!60!black}{\bfseries PASS} \\
  Commande introuvable           & GET    & 404 Not Found  & \textcolor{green!60!black}{\bfseries PASS} \\
  Supprimer sans droits          & DELETE & 403 Forbidden  & \textcolor{orange}{\bfseries WARN} \\
  \todo{Scénario 6}              & \todo{} & \todo{}       & \todo{} \\
  \bottomrule
\end{tabular}
\end{center}

%──────────────────────────────────────────────────────────
\section{Tests end-to-end (E2E)}
\label{sec:tests:e2e}
%──────────────────────────────────────────────────────────

\todo{Décrire les tests E2E : scénarios testés (parcours utilisateur
complets), outil (Cypress, Selenium, Playwright), rapport d'exécution.}

\placeholder{Rapport d'exécution Cypress / Selenium — figures/e2e\_report.png}

%──────────────────────────────────────────────────────────
\section{Tests de performance}
\label{sec:tests:performance}
%──────────────────────────────────────────────────────────

\todo{Si des tests de charge ont été réalisés (JMeter, k6, Gatling),
présenter la configuration et les résultats.}

\begin{center}
\begin{tabular}{>{\sffamily\bfseries}l r r r r}
  \toprule
  \textbf{Endpoint}    & \textbf{VU} & \textbf{Temps moyen (ms)}
                        & \textbf{P95 (ms)} & \textbf{Erreurs} \\
  \midrule
  GET /api/orders      & \todo{50}  & \todo{85}   & \todo{145}  & \todo{0\%} \\
  POST /api/orders     & \todo{50}  & \todo{195}  & \todo{320}  & \todo{0\%} \\
  GET /api/dashboard   & \todo{50}  & \todo{250}  & \todo{410}  & \todo{0,2\%} \\
  \todo{Endpoint 4}    & \todo{}    & \todo{}     & \todo{}     & \todo{} \\
  \bottomrule
\end{tabular}
\end{center}

%──────────────────────────────────────────────────────────
\section{Synthèse des résultats}
\label{sec:tests:synthese}
%──────────────────────────────────────────────────────────

\begin{center}
\begin{tabular}{>{\sffamily\bfseries}l l c}
  \toprule
  \textbf{KPI}                    & \textbf{Objectif}        & \textbf{Résultat} \\
  \midrule
  Couverture de tests             & > 75\%                   & \textcolor{green!60!black}{\bfseries\todo{85,3\%}} \\
  Temps de réponse P95            & < 400 ms                 & \textcolor{green!60!black}{\bfseries\todo{320 ms}} \\
  Taux d'erreur en charge         & < 1\%                    & \textcolor{green!60!black}{\bfseries\todo{0,2\%}} \\
  User stories livrées            & 14/14                    & \textcolor{orange}{\bfseries\todo{12/14}} \\
  Incidents en recette            & 0 bloquant               & \textcolor{green!60!black}{\bfseries\todo{0}} \\
  \todo{KPI métier}               & \todo{cible}             & \todo{valeur} \\
  \bottomrule
\end{tabular}
\end{center}

\clearpage

%==========================================================
%  CHAPITRE 09 — CONCLUSION GÉNÉRALE ET PERSPECTIVES
%==========================================================
\chapter{Conclusion générale et perspectives}
\label{chap:conclusion}

%──────────────────────────────────────────────────────────
\section{Bilan du projet}
\label{sec:ccl:bilan}
%──────────────────────────────────────────────────────────

Ce stage de \todo{durée} au sein de \textbf{\CompanyName} avait pour
objectif principal \todo{rappeler l'objectif central en une phrase}.

Au terme de cette expérience, les principaux résultats obtenus sont :

\begin{itemize}
  \item \todo{Résultat 1 : livrable principal et son impact mesurable}
  \item \todo{Résultat 2 : amélioration quantifiée (temps, coût, qualité)}
  \item \todo{Résultat 3 : contribution à l'équipe ou au système existant}
  \item \todo{Résultat 4 : apprentissage ou montée en compétence notable}
\end{itemize}

\todo{Rédiger un paragraphe de synthèse (5--8 phrases) qui évalue
l'atteinte des objectifs fixés en introduction, en soulignant les
écarts éventuels et leurs causes.}

%──────────────────────────────────────────────────────────
\section{Apports personnels}
\label{sec:ccl:apports}
%──────────────────────────────────────────────────────────

\subsection{Compétences techniques acquises}

\todo{Dresser la liste des compétences techniques nouvelles ou consolidées :
langages, frameworks, outils, pratiques de développement (TDD, CI/CD,
clean code, etc.).}

\begin{itemize}
  \item \textbf{Maîtrise approfondie de} \tech{\todo{framework / langage}}
  \item \textbf{Découverte et mise en pratique de} \tech{\todo{outil / pratique}}
  \item \textbf{Compréhension de} \todo{concept architectural ou méthodologique}
\end{itemize}

\subsection{Compétences transversales}

\todo{Mentionner les soft-skills développées : communication, gestion du
temps, travail en équipe, autonomie, animation de réunions, etc.}

\begin{itemize}
  \item \textbf{Communication :} participation aux daily standups et
        présentations d'avancement devant l'équipe.
  \item \textbf{Gestion de projet :} utilisation de \todo{Jira/Trello}
        pour le suivi des tâches, respect des délais de sprint.
  \item \textbf{Travail en équipe :} collaboration en mode pair-programming
        avec \todo{collègue}, revues de code croisées.
  \item \todo{Autre compétence transversale}
\end{itemize}

%──────────────────────────────────────────────────────────
\section{Difficultés rencontrées et leçons apprises}
\label{sec:ccl:difficultes}
%──────────────────────────────────────────────────────────

\todo{Synthétiser les principales difficultés (techniques, organisationnelles,
humaines) et les solutions trouvées. Se montrer critique et honnête :
ce qui a bien fonctionné, ce qui aurait pu être fait différemment.}

\begin{itemize}
  \item \textbf{Difficulté technique :} \todo{description + enseignement tiré}
  \item \textbf{Difficulté méthodologique :} \todo{description + enseignement tiré}
  \item \textbf{Difficulté organisationnelle :} \todo{description + enseignement tiré}
\end{itemize}

%──────────────────────────────────────────────────────────
\section{Perspectives et améliorations futures}
\label{sec:ccl:perspectives}
%──────────────────────────────────────────────────────────

Bien que les objectifs principaux aient été atteints, plusieurs pistes
d'amélioration et d'évolution sont envisagées à court et moyen terme :

\begin{enumerate}
  \item \textbf{Court terme (0--3 mois) :}
    \begin{itemize}
      \item \todo{Fonctionnalité non implémentée par manque de temps}
      \item \todo{Amélioration de performances identifiée en tests de charge}
      \item \todo{Correction de bugs mineurs détectés en recette}
    \end{itemize}

  \item \textbf{Moyen terme (3--12 mois) :}
    \begin{itemize}
      \item \todo{Extension du périmètre fonctionnel — nouvelle feature}
      \item \todo{Migration vers une architecture plus scalable}
      \item \todo{Intégration avec un système tiers}
    \end{itemize}

  \item \textbf{Long terme (> 1 an) :}
    \begin{itemize}
      \item \todo{Vision stratégique : digitalisation complète du processus X}
      \item \todo{Adoption de l'IA / ML pour automatiser la prise de décision}
      \item \todo{Autre perspective}
    \end{itemize}
\end{enumerate}

%──────────────────────────────────────────────────────────
\section{Mot de clôture}
\label{sec:ccl:mot}
%──────────────────────────────────────────────────────────

\begin{tcolorbox}[
  enhanced,
  colback=LightBg, colframe=ChapBlue,
  fontupper=\itshape,
  width=\textwidth
]
\todo{Rédiger un court paragraphe conclusif personnel (3--5 phrases)
qui exprime votre ressenti global, l'impact de ce stage sur votre
projet professionnel, et vos ambitions pour la suite de votre parcours.
Exemple : « Ce stage a été une expérience fondatrice qui a confirmé
mon attrait pour [domaine]. Il m'a permis de \ldots\ et m'a donné
envie de \ldots »}
\end{tcolorbox}

\clearpage


%──────────────────────────────────────────────────────────
%  BIBLIOGRAPHIE
%──────────────────────────────────────────────────────────
\clearpage
\addcontentsline{toc}{chapter}{Bibliographie}
\bibliographystyle{IEEEtran}   % ou plainnat pour APA
\bibliography{bibliography/references}

%──────────────────────────────────────────────────────────
%  ANNEXES
%──────────────────────────────────────────────────────────
\appendix
%==========================================================
%  CHAPITRE 10 — ANNEXES
%==========================================================
\chapter{Annexes}
\label{chap:annexes}

%──────────────────────────────────────────────────────────
\section{Annexe A — Cahier des charges détaillé}
\label{ann:cdc}
%──────────────────────────────────────────────────────────

\todo{Insérer ici le cahier des charges complet fourni par le maître
de stage, ou y faire référence si trop volumineux.}

\placeholder{Cahier des charges complet (document PDF joint ou extrait
— à insérer ici si applicable)}

%──────────────────────────────────────────────────────────
\section{Annexe B — Glossaire}
\label{ann:glossaire}
%──────────────────────────────────────────────────────────

\begin{description}[leftmargin=3cm, style=sameline]
  \item[\textbf{API}]  \textit{Application Programming Interface} —
        Interface permettant à deux logiciels de communiquer.
  \item[\textbf{CI/CD}] \textit{Continuous Integration / Continuous Delivery} —
        Pratiques DevOps automatisant les phases de build, test et déploiement.
  \item[\textbf{CRUD}] \textit{Create, Read, Update, Delete} —
        Les quatre opérations de base sur une base de données.
  \item[\textbf{DTO}]  \textit{Data Transfer Object} —
        Objet utilisé pour transporter des données entre les couches
        d'une application.
  \item[\textbf{JWT}]  \textit{JSON Web Token} —
        Standard ouvert pour l'authentification sans état.
  \item[\textbf{MVC}]  \textit{Model-View-Controller} —
        Patron architectural séparant la logique métier, la vue et le contrôleur.
  \item[\textbf{ORM}]  \textit{Object-Relational Mapping} —
        Technique de correspondance entre objets et tables relationnelles.
  \item[\textbf{REST}] \textit{Representational State Transfer} —
        Style architectural pour les services web.
  \item[\textbf{Scrum}] Cadre Agile de développement itératif basé sur des sprints.
  \item[\textbf{SLA}]  \textit{Service Level Agreement} —
        Accord définissant le niveau de service attendu.
  \item[\textbf{UML}]  \textit{Unified Modeling Language} —
        Langage de modélisation graphique standardisé.
  \item[\todo{Terme}]  \todo{Définition}
\end{description}

%──────────────────────────────────────────────────────────
\section{Annexe C — Documentation de l'API}
\label{ann:api}
%──────────────────────────────────────────────────────────

\todo{Insérer ici la documentation Swagger / OpenAPI des endpoints
développés, ou un tableau récapitulatif des routes.}

\begin{center}
\begin{tabular}{>{\ttfamily\footnotesize}l >{\sffamily\footnotesize\bfseries}c
                >{\footnotesize}l >{\footnotesize}l}
  \toprule
  \textbf{Route}              & \textbf{Méthode} & \textbf{Description}
                                                  & \textbf{Auth} \\
  \midrule
  /api/v1/auth/login          & POST   & Authentification   & Non \\
  /api/v1/auth/refresh        & POST   & Rafraîchir token   & Oui \\
  /api/v1/\todo{resource}     & GET    & Lister resources   & Oui \\
  /api/v1/\todo{resource}     & POST   & Créer resource     & Oui \\
  /api/v1/\todo{resource}/\{id\}& GET  & Lire resource      & Oui \\
  /api/v1/\todo{resource}/\{id\}& PUT  & Modifier resource  & Oui \\
  /api/v1/\todo{resource}/\{id\}& DELETE& Supprimer resource & Admin \\
  \todo{Autres routes}        & \todo{}& \todo{}           & \todo{} \\
  \bottomrule
\end{tabular}
\end{center}

%──────────────────────────────────────────────────────────
\section{Annexe D — Guide d'installation}
\label{ann:install}
%──────────────────────────────────────────────────────────

\subsection{Prérequis}

\begin{itemize}
  \item \todo{Java 17+ / Python 3.11+ / Node.js 18+}
  \item \todo{Docker 24.0+} et Docker Compose v2
  \item \todo{Git 2.40+}
  \item \todo{Autre prérequis}
\end{itemize}

\subsection{Étapes d'installation}

\begin{lstlisting}[
  language={},
  caption={Commandes d'installation et de lancement},
  label={lst:install}
]
# 1. Cloner le depot
git clone https://github.com/[A COMPLETER]/[projet].git
cd [projet]

# 2. Copier la configuration
cp .env.example .env
# Editer .env avec vos valeurs (DB, secrets, etc.)

# 3. Lancer avec Docker Compose
docker compose up --build -d

# 4. Initialiser la base de donnees (si applicable)
docker compose exec backend [commande de migration]

# 5. Acceder a l'application
# Frontend : http://localhost:3000
# API :      http://localhost:8080/api/v1
# Swagger :  http://localhost:8080/swagger-ui
\end{lstlisting}

%──────────────────────────────────────────────────────────
\section{Annexe E — Planning Gantt}
\label{ann:gantt}
%──────────────────────────────────────────────────────────

\placeholder{Diagramme de Gantt du projet — figures/gantt.png\\
(à générer avec MS Project, Excel, GanttProject ou pgfgantt)}

\todo{Commenter le Gantt : jalons importants, tâches critiques,
marges, décalages par rapport au plan initial.}


\end{document}
